\chapter{FIFO memory}\label{fifoMem}

\minitoc

The FIFO memory, or the {\em queue memory} is used to interconnect
subsystems working logical, or both logical and electrical,
asynchronously.x

\begin{definition}
FIFO memory implements a data structure which consists in a string
of maximum $2^{m}$ $n$-bit recordings accessed for write and read,
at its two ends. {\bf Full} and {\bf empty} signals are provided
indicating the write operation or the read operation are not
allowed. $\diamond$
\end{definition}

\placedrawing[h]{figures/G1_01_fifo_memory.lp  f66.lp}{{\bf FIFO memory.} {\small Two
pointers, evolving in the same direction, and a two-port RAM
implement a LIFO (queue) memory. The limit flags are computed
combinational from the addresses used to write and to read the
memory.}}{fifo}

A FIFO is considered synchronous if both {\em read} and {\em
write} signals are synchronized with the same clock signal. If the
two commands, {\em read} and {\em write}, are synchronized with
different clock signals, then the FIFO memory is called
asynchronous.

In Figure \ref{fifo} is presented a solution for the synchronous
version, where:
\begin{itemize}
    \item {\bf RAM} is a $2^{n}$ $m$-bit words two-port asynchronous random access memory, one
    port for write to the address {\tt w\_addr} and another for
    read form the address {\tt r\_addr}
    \item {\bf write counter} is an $(n+1)$-bit resetable counter incremented
    each time a write is executed; its output is {\tt w\_addr[n:0]}, initially
    it is reset
    \item {\bf read counter} is an $(n+1)$-bit resetable counter incremented
    each time a read is executed; its output is {\tt r\_addr[n:0]}, initially it is reset
    \item {\bf eq} is a comparator activating its output when
    the least significant $n$ bits of the two counters are
    identical.
\end{itemize}

FIFO works like a circular memory addressed by two pointers ({\tt
w\_addr[n-1:0]} and {\tt r\_addr[n-1:0]}) running on the same
direction. If the write pointer {\em after} a write operation
becomes equal with the read pointer, then the memory is  full and
the {\tt full} signal is 1. If the read pointer {\em after} a read
operation becomes equal with the write pointer, then the memory is
empty and the {\tt empty} signal is 1. The $n+1$-th bit in each
counter is used to differentiate between empty and full when {\tt
w\_addr[n-1:0]} and {\tt r\_addr[n-1:0]} are the same. If {\tt
w\_addr[n]} and {\tt r\_addr[n]} are different, then {\tt
w\_addr[n-1:0]} = {\tt r\_addr[n-1:0]} means full, else it means
empty.

The circuit used to compare the two addresses is a combinational
one. Therefore, its output has a hazardous behavior which affects
the outputs {\tt full} and {\tt empty}. These two outputs must be
used carefully in designing the system which includes this FIFO
memory. The problem can be managed because the system works in the
same clock domain ({\tt clock} is the same for both ends of FIFO
and for the entire system). We call this kind of FIFO {\em
synchronous FIFO}.

\begin{simulation} A Verilog synthesisable description of a synchronous
FIFO follows:

%\pagebreak

{\small
{\em
\begin{shaded*}
\begin{lstlisting}[language=Verilog, caption= , morekeywords={loop, repeat, doInParallel, for}]
/******************************************************************
File name:      simple_fifo.v
Circuit name:   Synchronous First-In First_Out memory
Description:    behavioral description of a synchronous FIFO
******************************************************************/
 module simple_fifo(output  [31:0]  out     ,
                    output          empty   ,
                    output          full    ,
                    input   [31:0]  in      ,
                    input           write   ,
                    input           read    ,
                    input           reset   ,
                    input           clock   );
    wire    [9:0]   write_addr, read_addr;

    counter write_counter(  .out        (write_addr ),
                            .reset      (reset      ),
                            .count_up   (write      ),
                            .clock      (clock      )),
            read_counter(   .out        (read_addr  ),
                            .reset      (reset      ),
                            .count_up   (read       ),
                            .clock      (clock      ));

    dual_ram memory(.out        (out            ),
                    .in         (in             ),
                    .read_addr  (read_addr[8:0] ),
                    .write_addr (write_addr[8:0]),
                    .we         (write          ),
                    .clock      (clock          ));

    assign  eq      = read_addr[8:0] == write_addr[8:0] ,
            phase   = ~(read_addr[9] == write_addr[9])  ,
            empty   = eq & phase   ,
            full    = eq & ~phase    ;
 endmodule
\end{lstlisting}
\end{shaded*}
}
}

%\pagebreak

{\small
{\em
\begin{shaded*}
\begin{lstlisting}[language=Verilog, caption= , morekeywords={loop, repeat, doInParallel, for}]
/******************************************************************
File name:      counter.v
Circuit name:   Counter
Description:    behavioral description of the counters used to 
                implement the asynchronous FIFO
******************************************************************/
 module counter(output  reg [9:0]   out     ,
                input               reset   ,
                input               count_up,
                input               clock   );

    always @(posedge clock) if (reset)          out <= 0;
                             else if (count_up) out <= out + 1;
 endmodule
\end{lstlisting}
\end{shaded*}
}
}

\pagebreak

{\small
{\em
\begin{shaded*}
\begin{lstlisting}[language=Verilog, caption= , morekeywords={loop, repeat, doInParallel, for}]
/******************************************************************
File name:      dual_ram.v
Circuit name:   Dual-Port RAM
Description:    behavioral description of the dual-port RAM used 
                to implement the synchronous FIFO
******************************************************************/
 module dual_ram(   output  [31:0]  out         ,
                    input       [31:0]  in          ,
                    input       [8:0]   read_addr   ,
                    input       [8:0]   write_addr  ,
                    input               we          ,
                    input               clock       );
    reg     [63:0]  mem[511:0];
    assign   out = mem[read_addr]   ;
    always @(posedge clock) if (we) mem[write_addr] <= in   ;
 endmodule
\end{lstlisting}
\end{shaded*}
}
}
 $\diamond$
\end{simulation}

An {\em asynchronous FIFO} uses two independent clocks, one for
{\tt write counter} and another for {\tt read counter}. This type
of FIFO is used to interconnect subsystems working in different
clock domains. The previously described circuit is unable to work
as an asynchronous FIFO. The signals {\tt empty} and {\tt full}
are meaningless, being generated in two clock domains. Indeed,
{\tt write counter} and {\tt read counter} are triggered by
different clocks generating the signal {\tt eq} with hazardous
transitions related to two different clocks: {\em write clock} and
{\em read clock}. This signal can not be used neither in the
system working with {\em write clock} nor in the system working
with {\em read clock}. {\em Read clock} is unable to avoid the
hazard generated by {\em write clock}, and {\em write clock} is
unable to avoid the hazard generated by {\em read clock}. Special
tricks must be used.
